\section{合成路线与相控制}

\subsection{固相与溶液前驱体}

常规固相路线以 BaCO$_3$/BaO、ZnO（或 ZnCO$_3$）和 SiO$_2$ 为原料，经混合、预烧、研磨和高温保温得到 BaZn$_2$Si$_2$O$_7$；Co、Mg 等替位通常沿用同一氧化物/碳酸盐反应路径，并通过中间研磨降低扩散距离 \cite{kerstan2013bazn2si2o7,thieme2016new,thieme2016cheminform}。燃烧辅助方法可缩短前驱体混合尺度，但 Eu$^{2+}$ 发光研究显示，最终相组成仍由后续煅烧温度和气氛决定 \cite{yao2010synthesis,yao2010synthesisx}。溶胶—凝胶路线把 Ba、Sr、Zn 和 Si 在分子尺度预混合，适合制备细粉和致密陶瓷；Kracker 等将其用于高、低温多晶型的可控获得 \cite{kracker2016sol}。

\subsection{玻璃析晶路线}

BaO/SrO/ZnO/SiO$_2$ 玻璃可通过成核—晶化热处理析出 Ba$_{1-x}$Sr$_x$Zn$_2$Si$_2$O$_7$。原位显微和高温衍射显示，Pt、Sb$_2$O$_3$、Nb$_2$O$_5$、Ta$_2$O$_5$、P$_2$O$_5$ 等成核剂会改变表面与体析晶的竞争关系 \cite{bocker2012in,vladislavova2017heterogeneous,vladislavova2019crystallization,heitmann2019surface,thieme2017on,vladislavova2018nucleation}。ZrO$_2$/TiO$_2$ 和 Al$_2$O$_3$ 添加剂则同时影响粘度、成核密度及最终晶粒尺寸 \cite{vladislavova2018crystallisation,thieme2017effect}。对于普通 BaZn$_2$Si$_2$O$_7$，玻璃路线的优势是成分均匀和晶粒可调，代价是残余玻璃相与微裂纹会掩盖本征膨胀 \cite{thieme2016thermal,thieme2015high}。

\begin{table}[t]
\centering
\caption{已报道路线的条件—结果对照。}
\begin{tabular}{P{0.19\linewidth}P{0.24\linewidth}P{0.23\linewidth}P{0.24\linewidth}}
\toprule
路线 & 前驱体/关键操作 & 主相判据 & 局限 \\
\midrule
混合氧化物固相 & 氧化物或碳酸盐；多次研磨和高温保温 & 粉末 XRD 与热膨胀一致 & 扩散受限、易残留 BaZnSiO$_4$ 等近邻相 \cite{kerstan2013bazn2si2o7} \\
溶胶—凝胶/燃烧 & 分子级混合；凝胶干燥或燃烧后煅烧 & 粒径小、可得到多晶型粉体 & 有机残基和气氛影响价态 \cite{kracker2016sol,yao2010synthesis} \\
玻璃析晶 & BaO/SrO/ZnO/SiO$_2$ 玻璃；成核+晶化 & 原位 XRD/显微观察晶化前沿 & 残余玻璃相、微裂纹、成核剂引入杂相 \cite{bocker2012in,vladislavova2017heterogeneous} \\
\bottomrule
\end{tabular}
\end{table}

\subsection{相图与热力学数据}

现有文献对 BaO--ZnO 二元关系有相平衡评估，但未检索到可直接用于 BaO--ZnO--SiO$_2$ 三元体系的完整 CALPHAD 数据库。因而，组成点是否落在 BaZn$_2$Si$_2$O$_7$ 初晶区、是否经过共晶液相，只能结合实验 XRD 和热分析判断。对负热膨胀硅酸盐的理论—实验工作表明，声子自由能和温度依赖的相稳定性不可由室温结构单独推断 \cite{erlebach2021phase}。
